\documentclass[11pt, a4paper]{article}

% Packages
\usepackage[utf8]{inputenc}
\usepackage[T1]{fontenc}
\usepackage{lmodern}
\usepackage[margin=0.75in]{geometry}
\usepackage{xcolor}
\usepackage{amsmath}
\usepackage{amssymb}
\usepackage{titlesec}
\usepackage{enumitem}
\usepackage{hyperref}
\usepackage{tikz}
\usepackage{tcolorbox}
\usepackage{multicol}
\usepackage{tabularx}
\usepackage{array}
\usepackage{parskip}
\usepackage{fancyhdr}

% === Grayscale Color Scheme ===
\definecolor{primary}{RGB}{60, 60, 60}
\definecolor{secondary}{RGB}{80, 80, 80}
\definecolor{accent}{RGB}{240, 240, 240}
\definecolor{textgray}{RGB}{70, 70, 70}
\definecolor{lightgray}{RGB}{245, 245, 245}

% Hyperlink setup
\hypersetup{
    colorlinks=true,
    linkcolor=primary,
    urlcolor=primary,
    pdftitle={Minje Park - CV},
    pdfauthor={Minje Park}
}

% Custom section formatting
\titleformat{\section}
    {\color{primary}\Large\bfseries\scshape}
    {}{0em}
    {}
    [\color{primary}\titlerule]

\titlespacing{\section}{0pt}{20pt}{8pt}

% Subsection formatting
\titleformat{\subsection}
    {\color{secondary}\large\bfseries}
    {}{0em}
    {}

% Page style
\newcommand{\cvrulewidth}{1.5pt}
\pagestyle{fancy}
\fancyhf{}
\fancyhead[R]{\small\color{secondary} Minje Park - Curriculum Vitae}
\renewcommand{\headrulewidth}{\cvrulewidth}
\renewcommand{\footrulewidth}{\cvrulewidth}
\lfoot{\color{textgray}\small\textit{Last updated: March 2026}}
\cfoot{\small - \thepage\ -}

% Custom commands
\newcommand{\cvtag}[1]{%
    \tikz[baseline=(tag.base)]\node[draw=primary, fill=accent, rounded corners=3pt, inner sep=3pt, text=primary] (tag) {\small #1};%
}

\newcommand{\experienceitem}[4]{
    \begin{tcolorbox}[
        colback=lightgray,
        colframe=primary,
        boxrule=0pt,
        leftrule=3pt,
        arc=0pt,
        left=5pt,
        right=5pt,
        top=5pt,
        bottom=5pt
    ]
    \textbf{\large #1} \hfill {\color{primary}#2}
    
    \textit{#3}
    
    #4
    \end{tcolorbox}
    \vspace{3pt}
}

\newcommand{\conferenceitem}[3]{%
    \noindent
    \parbox[t]{0.15\textwidth}{{\color{primary}\textbf{#1}}}%
    \hfill
    \parbox[t]{0.83\textwidth}{\textbf{#2}\\[2pt]\textit{#3}}%
    \vspace{8pt}\par
}

\newcommand{\awarditem}[3]{
    \begin{tcolorbox}[
        colback=white,
        colframe=primary,
        boxrule=0pt,
        leftrule=3pt,
        arc=0pt,
        left=5pt,
        right=5pt,
        top=4pt,
        bottom=4pt
    ]
    \textbf{#1} \\[2pt]
    #2 \\[2pt]
    {\small\textit{#3}}
    \end{tcolorbox}
    \vspace{3pt}
}

\begin{document}

% === Header ===
\vspace{-6pt}

{\Huge\bfseries\color{primary} MINJE PARK}

\vspace{4pt}
\noindent
\begin{tabularx}{\textwidth}{@{}>{\bfseries}p{0.16\textwidth}X@{}}
Email &
\href{mailto:pmj03400@naver.com}{pmj03400@naver.com} \\
Homepage &
\href{https://pmj0324.github.io/}{https://pmj0324.github.io/} \\
GitHub &
\href{https://github.com/pmj0324}{github.com/pmj0324} \\
Scholar &
\href{https://scholar.google.com/citations?user=oC5mE0AAAAAJ}{Google Scholar (Minje Park)} \\
\end{tabularx}

\vspace{2pt}

% === Education ===
\section{Education}

\noindent
\textbf{Sungkyunkwan University} \hfill \textit{Feb 2021 -- Feb 2027 (Expected)} \\[2pt]
B.S. in Physics, Double Major in Electronic \& Electrical Engineering

% === Research Experience ===
\section{Research Experience}

\experienceitem{GENESIS: Multi-Field Cosmological Simulation and Inference}{May 2025 -- Present}{Independent Research}{
    \begin{itemize}[leftmargin=*, itemsep=2pt]
        \item Designed GENESIS (GENerative AI for thE Simulation and Inference of cosmological fieldS), a flow-matching framework aimed at two related tasks: conditional cosmological field generation and training-free parameter inference.
        \[
        \frac{d\mathbf{x}_t}{dt} = \mathbf{v}_\theta(\mathbf{x}_t,t,\mathbf{c}), \qquad
        \mathcal{L}_{\mathrm{FM}} = \mathbb{E}\!\left[\left\|\mathbf{v}_\theta(\mathbf{x}_t,t,\mathbf{c}) - \mathbf{u}_t\right\|^2\right]
        \]
        \item The learned probability-flow vector field is integrated to generate dark matter, gas, and temperature fields conditioned on CAMELS parameters, while the same flow is also used for likelihood-based inference without additional training.
        \[
        \log p_\theta(\mathbf{x}_1 \mid \mathbf{c}) = \log p(\mathbf{x}_0) - \int_0^1 \nabla \cdot \mathbf{v}_\theta(\mathbf{x}_t,t,\mathbf{c})\,dt
        \]
        Studying maximum-likelihood estimation of cosmological parameters $\mathbf{c}$ from the same learned flow.
    \end{itemize}
}

\experienceitem{CNN and Vision Transformer for the IceCube OPTICUS Camera System}{Aug 2023 -- Sep 2025}{Supervisor: Carsten Rott, Department of Physics, Sungkyunkwan University}{
    \begin{itemize}[leftmargin=*, itemsep=2pt]
        \item Developed CNN and Vision Transformer models for the IceCube Upgrade Camera System; results published as a poster and proceedings paper at ICRC 2025 (arXiv:2507.18525).
    \end{itemize}
}

\experienceitem{Generative Modeling for IceCube and SWGO}{May 2024 -- Present}{Supervisor: Chang Dong Rho, Department of Physics, Sungkyunkwan University}{
    \begin{itemize}[leftmargin=*, itemsep=2pt]
        \item Developed diffusion models for SWGO detector-signal synthesis.
        \item Built JANUS, a flow-matching framework for conditional IceCube detector-signal generation and event-level inference.
    \end{itemize}
}

\experienceitem{Transformer-Based Reconstruction for SWGO}{Jun 2024 -- Present}{Supervisor: Ian Watson, Department of Physics, University of Seoul}{
    \begin{itemize}[leftmargin=*, itemsep=2pt]
        \item Investigated position embedding strategies for transformer-based primary-particle reconstruction in SWGO.
        \item Studied flow-matching models for SWGO detector-signal synthesis from simulation data.
    \end{itemize}
}

% === Publications ===
\section{Publications}

\newcommand{\publication}[3]{
    \vspace{4pt}
    \noindent
    \textbf{#1} \\
    \textit{#2} \\
    {\small #3}
    \vspace{6pt}
}

\publication{Performance Study of the IceCube Upgrade Camera System}
{C. Rott, \textbf{\textit{M. Park}}, M. Jansson, G. Iverson, S. Choi (for the IceCube Collaboration)}
{Proceedings of Science (PoS), ICRC2025; \href{https://arxiv.org/abs/2507.18525}{arXiv:2507.18525}}

% === Conference Presentations ===
\section{Conference Presentations}
\conferenceitem{Apr 2026}{[Talk, scheduled] Generative AI for IceCube Observatory}{102nd Spring Meeting of the Korean Physical Society, Daejeon, Republic of Korea}

\conferenceitem{Apr 2026}{[Poster, scheduled] Mixture Density Networks for Fast Detector Response Simulation in SWGO}{102nd Spring Meeting of the Korean Physical Society, Daejeon, Republic of Korea}

\conferenceitem{Apr 2026}{[Poster, scheduled] Flow Matching Model for Generating Cosmic Web Images}{102nd Spring Meeting of the Korean Physical Society, Daejeon, Republic of Korea}

\conferenceitem{Feb 2026}{[Poster] Investigation of IGZO-Based 3D Vertical Channel Transistors and Physics-Informed Neural Network Modeling for Next-Generation DRAM}{The 33rd Korean Conference on Semiconductors, Jeongseon, Republic of Korea}

\conferenceitem{Oct 2025}{[Poster] Generative AI for Cosmological Simulation and Manifold Representation}{KPS Undergraduate Physics Research Competition, 102nd Fall Meeting of the Korean Physical Society, Gwangju, Republic of Korea}

\conferenceitem{Jun 2025}{[Poster] Diffusion and Flow Matching Model for IceCube Neutrino Event Simulation and Reconstruction}{4th K-Neutrino Symposium, Seoul, Republic of Korea}

\conferenceitem{Apr 2025}{[Poster] Modeling the Relationship between PMT Response and Air-Shower Particles for SWGO}{101st Spring Meeting of the Korean Physical Society, Daejeon, Republic of Korea}

\conferenceitem{Oct 2024}{[Poster] Development of a Deep Learning Model for Analyzing Antarctic Ice Optical Properties}{KPS Undergraduate Physics Research Competition, 101st Fall Meeting of the Korean Physical Society, Yeosu, Republic of Korea}

\conferenceitem{Sep 2024}{[Talk] Deep Learning for Ice-Property Measurement Using Simulation Images}{Fall IceCube Collaboration Meeting, Madison, Wisconsin, USA}

\conferenceitem{Jul 2024}{[Talk] Ice-Property Measurement Using a Camera System in the IceCube Upgrade}{3rd K-Neutrino Symposium, Gwangju, Republic of Korea}

\conferenceitem{Apr 2024}{[Poster] Study of Optical Properties in Antarctic Ice Using Simulated Images}{100th Spring Meeting of the Korean Physical Society, Daejeon, Republic of Korea}

% === Honors & Awards ===
\section{Honors \& Awards}

\awarditem{Best Presentation Award}{4th K-Neutrino Symposium, 2025}{Awarded for poster presentation: ``Diffusion and Flow Matching Model for IceCube Neutrino Event Simulation and Reconstruction''}

\awarditem{Merit Award}{Korean Physical Society Undergraduate Physics Research Competition, 2024}{Recognized for research on CNN and Vision Transformer models for the IceCube Upgrade Camera System}

% === Skills ===
\section{Skills}

\begin{tabularx}{\textwidth}{@{}>{\bfseries}p{0.18\textwidth}X@{}}
Programming   & Python, NumPy, SciPy, Matplotlib \\[4pt]
Frameworks    & PyTorch, JAX \\[4pt]
Models        & Transformer, Vision Transformer (ViT), CNN, Diffusion Models, Flow Matching \\[4pt]
Tools         & Linux, HTCondor, Git, Bash/Shell, HDF5 \\
\end{tabularx}

% === Collaboration Memberships ===
\section{Collaboration Memberships}

\noindent
\textbf{IceCube Collaboration} \hfill \textit{Aug 2023 -- Present} \\
{\small A neutrino detector embedded in Antarctic ice at the South Pole}

\vspace{8pt}

\noindent
\textbf{SWGO Collaboration} \hfill \textit{Jun 2024 -- Present} \\
{\small A next-generation ground-based gamma-ray observatory planned for the high-altitude Andes}

\end{document}
